% 第 3 章 §3.1（印刷页 109–118）
\chapter{拓扑场论五讲}
\label{chap:03}

\noindent Constantin Teleman\footnote{本文是 2012 年春在巴塞罗那 Bellaterra 的 CRM 所作系列讲义的修订版。作者感谢主办机构的热情接待。本笔记的第一稿写于石溪 Simons 中心访问期间，其中的一些想法源自与 Dan Freed 的合作工作。此外，作者就许多富有启发的讨论感谢 Chris Douglas、Mike Hopkins、Jacob Lurie 与 Chris Schommer-Pries。工作完成期间的部分经费由 NSF 通过 DMS--1160328 号基金提供。}

\section{引言与例子}

拓扑量子场论（topological quantum field theory，TQFT）在物理学中作为真正量子场论的低能（零能量）部分出现，它表现出对时空大尺度拓扑的意外依赖。态空间的零能量部分不随时间演化——按定义，它被哈密顿量杀死——所以初看之下，它的物理似乎无趣。但这个论证没有考虑拓扑非平凡的时空。

在数学中，TQFT 是作为流形的某些新的拓扑或微分不变量的组织框架而出现的：这些不变量无法被代数拓扑的标准技术捕捉。（我们将会看到其中的原因。）

\subsection{定义}

Atiyah 与 Witten 的原始公理（见 [30]）受 Graeme Segal 的共形场论公理启发，把 $D$ 维 TQFT 定义为一个对称幺半函子
\[
Z\colon\bigl(\mathsf{Bord}^{\mathrm{or}}_D,\sqcup\bigr)\longrightarrow\bigl(\mathsf{Vect},\otimes\bigr).
\]
这里 $\mathsf{Bord}^{\mathrm{or}}_D$ 的对象是 $D-1$ 维紧的无边界定向流形，态射是定向的 $n$ 维配边，模去相对于边界的微分同胚。\footnote{这个颇为朴素的定义对简单应用已经够用。一个改进的定义把微分同胚群的拓扑考虑进来，得到在拓扑群胚中丰富的范畴，其态射是由配边与相对于边界的微分同胚构成的拓扑群胚。} 配边假定有入边界与出边界；但在定向的世界里，这一选择也可以通过比较每个分支上的边界定向与独立定向来指明。于是，有两个正端点的区间必然有一个入端与一个出端。

\begin{remark}
在一些精确的表述中，流形被嵌入一个非常大的欧氏空间（$\R^\infty$）中，而配边嵌入到额外的一个「时间」维度中。我们将忽略这一结构。
\end{remark}

配边范畴有对称幺半结构（一个结合且交换的乘法函子 $\sqcup\colon\mathsf{Bord}^{\mathrm{or}}_D\times\mathsf{Bord}^{\mathrm{or}}_D\to\mathsf{Bord}^{\mathrm{or}}_D$），由不交并定义。向量空间范畴有类似的乘法，由张量积定义。由于 $\mathsf{Vect}$ 的额外线性性质与 $\otimes$ 的双线性，我们称之为张量结构。

\subsection{例：有限群规范理论}

与有限群 $F$ 相联系的这一理论 $Z_F$，在任意维数都容易构造。缺点是，它只能探测流形的基本群。\footnote{背景中还潜藏着一个更精细的增强，即有限同伦型的规范理论，它探测更高的同伦信息，并与群代数的范畴化概念相关；见 §3.3。}

记 $\mathrm{Bun}_F(X)$ 为 $X$ 上主 $F$-丛的群胚。对一个闭 $D$ 维流形 $M$，$Z_F$ 赋值给它的，是 $M$ 上主 $F$-丛的同构类的个数，每一类按它的自同构群的阶加权（即除以该阶）。对连通的 $M$，这个数是 $\#\mathrm{Hom}(\pi_1(M),F)/\#F$。于是当 $D=1$ 时 $Z_F(S^1)=1$。

对一个闭的 $(D-1)$ 维流形 $N$，$Z_F(N)$ 表示 $\mathrm{Bun}_F(N)$ 上的函数空间。

现在应当容易猜到：$Z_F(M)\colon Z_F(\partial_-M)\to Z_F(\partial_+M)$ 是这样的线性映射——它联系 $F_-\to\partial_-M$ 与 $F_+\to\partial_+M$ 的矩阵元，计数 $M$ 上限制到这两个边界上指定丛的 $F$-丛。计数按在出边界 $\partial_+M$ 上为平凡的自同构加权。验证这给出一个 TQFT（即配边映射的复合对应到线性映射的复合）是一个练习，最好按下述注记展开。

\begin{remark}[对应图]
这个 TQFT 隐秘地由一个群胚的对应定义，我们将在 §3.2 遇到这种对应。即，群胚 $\mathrm{Bun}_F(M)$ 通过限制映射到 $\mathrm{Bun}_F(\partial_-M)$ 与 $\mathrm{Bun}_F(\partial_+M)$：
\begin{center}
\begin{tikzcd}[row sep=1.6em, column sep=2.6em]
& \mathrm{Bun}_F(M) \arrow[dl, "\pi_-"'] \arrow[dr, "\pi_+"] & \\
\mathrm{Bun}_F(\partial_-M) & & \mathrm{Bun}_F(\partial_+M).
\end{tikzcd}
\end{center}
把 $Z_F(\partial_\pm M)$ 等同于相应群胚的 $H^0(\cdot;\C)$；映射 $Z_F(M)$ 则是上同调上的「推–拉」复合 $(\pi_+)_*\circ\pi_-^*$。这里我们约定：沿群胚的推前 $(\pi_+)_*$ 对点求和，但每个点按它的自同构群的阶的倒数加权。
\end{remark}

\begin{remark}[扭转]
上述有限规范 TQFT 不使用定向，但它的一个扭转变体要用。取一个群上同调类 $\tau\in H^D(BF;\C^\times)$；它对每个主 $F$-丛 $P\to M$ 定义一个类 $\tau_P\in H^n(M;\C^\times)$。在最高维，我们定义 $Z^\tau_F(M)$ 为：在加权计数中，对每个 $P$ 再乘上复数 $\int_M\tau_P$。转到闭的 $(D-1)$ 维流形 $N$，稍加思考便知 $\tau$ 定义了 $N$ 上每个主丛的自同构群的一维特征标。（提示：用自同构在 $S^1\times N$ 上定义一个丛，再用 $\tau$ 为每个自同构提取一个复数。）在定义 $Z^\tau_F(N)$ 时，我们删去所有对应于自同构非平凡作用的主丛的项。换言之：$\tau$ 定义了 $\mathrm{Bun}_F(N)$ 上的一个平坦线丛 $\mathcal{O}(\tau)$，而 $Z^\tau_F(N)=H^0(\mathrm{Bun}_F(N;\mathcal{O}(\tau)))$。
\end{remark}

对应描述可推广到扭转情形，并给出数学物理中路径积分的一个雏形例子——构造量子场论的一种受人喜爱（尽管通常不严格）的技术。回忆在物理学中，联系于闭流形的配分函数是场空间上的积分 $\int_{\text{场 }\phi}\exp\{iS(\phi)\}\,D\phi$，其中 $S(\phi)$ 是「作用量」。在这里，场是到 $BF$ 的映射，即 $M$ 上的主丛；测度是自同构数的倒数；而（指数化）作用量的角色由 $\tau$ 扮演。

\begin{exercise}
对带边界的 $D$ 维流形，显式描述 $Z^\tau_F$ 的矩阵系数。
\end{exercise}

\subsection{初步分类，$D=1$}

我们约定：在图中时间从右向左流动，以匹配算子的复合。向量空间 $Z(+)$、$Z(-)$ 赋值给带两种定向的点。右弧 $\supset$ 给出一个态射 $Z(\supset)\colon\C=Z(\varnothing)\to Z(+)\otimes Z(-)$，左弧给出一个配对 $Z(+)\otimes Z(-)\to\C$。只在一个因子中作适当的复合，便知这两个映射在 $Z(+)$ 与 $Z(-)$ 之间建立完美对偶（必要时细节见 §3.4）；这特别迫使它们有限维：否则 $Z(+)$ 上的恒等映射就不会落在 $Z(+)\otimes Z(+)^\vee$ 中。于是 $Z$ 由 $Z(+)=V$ 描述；$Z(S^1)=\dim V$，而所有运算都涉及 $V\otimes V^\vee$ 中的标准展开与收缩张量。

值得注意的是，这个雏形例子已经包含 Lurie 配边假设的胚芽——该假设用数据 $Z(+)$ 来分类（全扩张）TQFT，而 $Z(+)$ 必须满足一个用对偶性表述的有限性假设。我们将在 §3.4 中阐明这一点。

\begin{remark}[无定向情形]
即使 $D=1$ 的情形，若放弃定向也会变得有趣。此时必须有 $Z(+)=Z(-)$，因为只有一个点；而在这个阶段，我们得到的是一个带非退化对称双线性型的向量空间。（练习：用图证明该型的对称性。）在更高维，我们选择定向流形而不是标架流形\footnote{这一区别在一维不可见。}，这一选择将对我们的理论施加强限制，它与复几何中的 Calabi--Yau 条件相关：典范丛的平凡化。
\end{remark}

\subsection{$D=2$ 与Frobenius 代数}

按我们给出的定义来分类 $D=2$ 的 TQFT 并不太难；这个结果流传已有一段时间，Abrams [1] 给出了严格的表述。

\begin{theorem}
一个二维定向 TQFT 等价于一个有限维向量空间 $A$ 上的交换Frobenius 代数结构，其中 $A$ 正是 $Z(S^1)$。
\end{theorem}

这里，一个Frobenius 代数（Frobenius algebra）是一个带单位的结合代数 $A$，连同一个迹 $\theta\colon A\to\C$，它通过 $a\times b\mapsto\theta(ab)$ 诱导向量空间 $A\times A\to\C$ 上的非退化配对。迹条件 $\theta(ab)=\theta(ba)$ 在非交换情形变得重要，并将在二维扩张 TQFT 中出现。交换Frobenius 代数的一个例子是闭定向流形的上同调，以杯积为乘法、以积分映射为迹。把杯积形变为量子上同调——在射影流形上计数全纯曲线——则与著名的 Gromov--Witten 理论相关。

各运算的几何表示是众所周知的：乘法由裤子表示，它映射 $Z(S^1)\otimes Z(S^1)\to Z(S^1)$；单位由带出边界的圆盘表示；迹由出圆盘表示。迹配对 $\theta(ab)$ 由一个「肘形」实现，即带两个入圆圈的弯曲柱面。由图形容易推出乘法的交换性与迹配对的非退化性（与一维情形中 $Z(\subset)$ 的非退化性相对应）；但定理的反方向需要多花一点功夫。下面是论证的概要：Frobenius 代数含有一个特殊元素，即欧拉类 $\alpha$，它是由带单个出圆圈的双环面输出的向量。于是，闭的亏格 $g$ 曲面算出 $\theta(\alpha^g)$，而带 $p$ 个入边界与 $q>0$ 个出边界的曲面算出如下映射 $A^{\otimes p}\to A^{\otimes q}$：$p$ 个输入的乘积乘以 $\alpha^g$，再余乘到输出中。请读者验证这个公式满足 Atiyah--Witten 定义 3.1.1。

代数半单的特殊情形值得注意。此时 $A=\bigoplus_i\C P_i$，其中投影算子满足 $P_iP_j=\delta_{ij}P_i$，而它们的迹 $\theta_i:=\theta(P_i)$ 必须是非零复数。这些数在差一个次序的意义下决定了Frobenius 代数的同构类。于是 $\alpha=\sum_i\theta_i^{-1}P_i$，且 $Z(\Sigma_g)=\sum_i\theta_i^{1-g}$。

每当某个与曲面相关的不变量能写成 $\sum_i\theta_i^{1-g}$ 的形式时，就应当怀疑它受一个二维 TQFT 支配。

\subsection{二维的有限群规范理论}

让我们把例 3.1.2 在两维中写开。圆周上主 $F$-丛的同构类是 $F$ 中的共轭类。其上的函数空间以特征标为自然基。由例 3.1.2 中的推–拉图，可以算出由裤子描述的乘法运算，得到特征标（视为群上的函数）的卷积。于是该代数是半单的，其投影算子由类函数 $(\#F)^{-1}\dim V\cdot\chi_V$ 给出，对各个不可约表示 $V$ 而言。代数的单位是原点处的 delta 函数，即投影算子之和 $(\#F)^{-1}\sum_V\chi_V\cdot\dim V$。由出圆盘定义的迹，是类函数在单位元处的取值——即圆周上唯一可扩张到圆盘的 $F$-丛的完整和乐——再按平凡丛的自同构数 $\#F$ 加权。于是投影算子的迹为 $\theta_V=\dim^2V/(\#F)^2$；比较亏格 $g$ 配分函数的两种算法，便得到恒等式
\[
(\#F)^{2g-2}\sum_V(\dim V)^{2-2g}
=\frac{\#\{u_1,u_2,\dots,u_{2g}\in F\mid [u_1,u_2]\cdot[u_3,u_4]\cdots[u_{2g-1},u_{2g}]=1\}}{\#F}.
\]
在亏格 $0$ 与 $1$，它们是有限群特征标论中的一些经典恒等式：$\sum(\dim V)^2=\#F$，以及不可约表示同构类的个数等于共轭类的个数。这个 TQFT 例子以代数恒等式的形式，早已为 Frobenius 所知。

\subsection{有限高阶群胚理论}

有限规范理论可以推广，以捕捉流形更多的同伦信息。为此，我们改用分类空间 $BF$ 来重新理解有限群 $F$。回忆它是可缩空间模去自由 $F$-作用的商。（一个熟知的例子是 $B\Z/2\cong\RP^\infty=S^\infty/(\Z/2)$。）这个 $BF$ 是 Eilenberg--MacLane 空间，即只有一个非平凡同伦群的空间：在这里 $\pi_1=F$。$M$ 上主丛的同构类是映射空间到 $BF$ 的连通分支 $[M,BF]$。套用物理学的说法，我们正在量子化取值于 $BF$ 的场，但「路径积分」数的是分支（带自同构，扭转情形下再按作用量加权）。

类似地，$(D-1)$ 维流形 $N$ 的态空间，是 $N$ 上取值于 $BF$ 的场的局部常值函数空间。扭转情形下，我们得到这个映射空间上的一个（平坦）线丛，并取其局部常值截面。

现在可以把 $BF$ 换成更一般的靶空间 $X$。所得理论的一个好的有限性概念，是 $X$ 的同伦群有限。此外，若 $\dim M\le D$，高于 $\pi_D$ 的同伦群不影响映射空间的分支集 $[M,X]$，所以我们取 $X$ 具有有限同伦群且在维数 $D$ 以上消失。我们这样构造一个 TQFT：令 $Z_X(N^{D-1})$ 为 $\mathrm{Map}(N;X)$ 上局部常值函数的向量空间，并令 $Z_X(M)$ 以某些权计数分支 $[M,X]$。所用的正确权稍难算出：对 $m\colon M\to X$，它是映射空间 $\mathrm{Map}(M,X)$ 的 $m$-分支的各同伦群阶的交错乘积。与规范理论一样，这个 TQFT 可以被一个「作用量」$\tau\in H^D(X;\C^\times)$ 扭转。

$\mathrm{Map}(M;X)$ 的同伦群可以用从 $H^p(M;\pi_{-q}X)$ 出发的谱序列计算。然而这个计算涉及关于 $M$ 与 $X$ 的更精细的信息，即 $M$ 上的杯积，以及 $X$ 的各同伦群如何层叠在一起（所谓的 Postnikov $k$-不变量）。所以上述 TQFT 比 $M$ 的同调与 $X$ 的同伦群更敏感。举一个简单的例子\footnote{尽管同伦群无限，为了停留在熟悉的空间上才这样取。把同伦群模 $n$ 塌缩即可得到有限的例子。}：回忆 $\pi_2(S^2)\cong\pi_3(S^2)\cong\Z$ 与 $\pi_4(S^2)\cong\pi_4(S^3)\cong\Z/2$。一个在不超过 $4$ 维与 $S^2$ 有相同同伦群的空间是 $\CP^\infty\times S^3$。现在，从 $M=\CP^2$ 到后者的映射很多：例如到因子 $\CP^\infty$ 的映射由拉回生成元 $\sigma\in H^2(\CP^\infty)$ 在 $H^2(\CP^2)$ 中的次数分类；等价地，由诱导映射 $\pi_2\CP^2\to\pi_2\CP^\infty$（两个群都等于 $\Z$）分类。这些有趣的映射没有一个来自映入球面 $S^2=\CP^1\subset\CP^\infty$ 的映射：障碍在于 $S^2$ 上面积形式 $\sigma$ 的杯平方为零，而 $\CP^2$ 上杯平方 $H^2\to H^4$ 不为零。所以面积形式 $\sigma\in H^2(S^2)$ 在任何映射 $\CP^2\to S^2$ 下都拉回为零。我们将在 §3.3 中以另一种面貌再次遇到这个二次障碍。

\subsection{二维 Yang--Mills 理论}

规范理论的最后一个变体——它将占据下一节——涉及紧规范群 $G$。刚刚讲述的关于计数主丛的故事需要解释。一种做法是把「主丛」解释为「平坦主丛」，即带平坦联络的丛。（对离散群，主丛有自然的平坦结构。）

如前，$S^1$ 上平坦丛的同构类的模空间等同于共轭类空间 $G/G^{\mathrm{ad}}$，所以联系于圆周的向量空间应当是群上的类函数空间，它同样带有特征标这一自然基。闭亏格 $g$ 曲面上平坦丛也有模空间 $F(\Sigma_g;G)$：它是紧的实解析空间，实维数为 $2(g-1)\cdot\dim(G)$（若 $G$ 单；一般情形须按中心的秩修正）。这个簇通常是奇异的，可以借助 $\pi_1(\Sigma_g)$ 描述为
\begin{equation}
\Phi(\Sigma,G)=\Bigl\{(u_1,u_2,\dots,u_{2g})\in G^{2g}\ \Bigm|\ [u_1,u_2]\cdots[u_{2g-1},u_{2g}]=1\Bigr\}\Big/ G,
\tag{3.1}\label{eq:3.1}
\end{equation}
其中群通过同时共轭作用。若要求换位子不取 $1$、而落在某个指定的共轭类 $c$ 中，就得到相关的簇 $\Phi_c$（维数不同）；对一般的 $c$，该簇是光滑的。这些簇与带边界曲面的 TQFT 有关。

具体地说，考虑带单个边界圆周的曲面，并把它取作出边界。由于 $\pi_1$ 自由，模空间是 $G^{2g}/G^{\mathrm{ad}}$。受注 3.1.2 中对应图的启发，我们研究「换位子乘积」映射 $G^{2g}\to G$（它对 $G$ 的共轭作用等变）；而对任意 $c$ 的簇 $\Phi_c$ 本质上就是这个映射的纤维。

「计数」丛不是一个选项，但另一种做法来自这样的观察：这些模空间有来自辛形式的自然体积元（见下一节）。若约定对这个体积积分而不是计数，我们几乎得到一个二维 TQFT，其亏格 $g$ 配分函数是空间 (3.1) 的辛体积。「几乎」是指：我们实际上得到的是无穷维向量空间 $Z(S^1)$，不过它带有自然基，即 $G$ 的不可约特征标。其背后的Frobenius 代数结构由 Witten [31] 算出，答案与有限群的情形惊人地相似：Witten 对 $G$ 的每个不可约特征标找到投影算子 $P_V$，其迹为 $\theta_V=\mathrm{vol}(G)^{-2}\dim^2V$。（$G$ 的黎曼体积用李代数 $\mathfrak{g}$ 上一个共轭不变度量计算；同一度量也用来定义辛形式。）由此 Witten 得到辛体积公式
\[
\int_{\Phi(\Sigma;G)}\exp(\omega)=\#Z(G)\cdot\mathrm{vol}(G)^{2g-2}\sum_V(\dim V)^{2-2g},
\]
若 $G$ 半单且亏格大于 $1$，这是一个收敛级数。

\begin{remark}
因子 $\#Z(G)$（$G$ 的中心的阶）似乎有些突兀。但事实上，中心元素定义了任何平坦丛的自同构，而且在亏格 $>2$ 时，一般情形下它们是仅有的自同构。所以 $\Phi$ 一般是以 $Z(G)$ 为稳定化子的轨形，积分时本应除以它。
\end{remark}

\subsection{TQFT 的变体：上同调场论}

一类相当著名的二维场论，通过记住曲面有微分同胚来增强我们讨论过的结构。它们就是镜像对称的 A-模型与 B-模型：前者在紧辛流形中计数（伪）全纯曲线，后者描述复流形上凝聚层上的运算，二者都产生取值于曲面丛的特征类中的不变量。

在最简单的表述中，我们保留联系于圆周的向量空间 $Z(S^1)$，但要求由带 $p$ 个输入与 $q$ 个输出的曲面诱导的线性映射 $Z(S^1)^{\otimes p}\to Z(S^1)^{\otimes q}$ 在族中「上同调地变化」。也就是说，若 $B$ 是曲面丛 $\Sigma\to B$ 的底，我们希望 $Z(\Sigma)$ 定义
\[
H^*\Bigl(B;\operatorname{Hom}\bigl(Z(S^1)^{\otimes p},Z(S^1)^{\otimes q}\bigr)\Bigr)
\]
中的一个类。$H^0$ 分量就是原来的线性映射；「底上的变化」不是指对点 $b\in B$ 的连续依赖，而是指与同调类的耦合。在拓扑语境中，这是我们想计算同伦不变量时唯一合理的定义。

这个概念有许多变体（见 [7, 29]），这里要标出的是 Kontsevich 与 Manin 在 [20] 中定义的上同调场论（cohomological field theory，CohFT），其中我们允许曲面出现节点，如同代数几何中的 Lefschetz 纤维化。这意味着在奇点附近局部地，曲线族由族 $(x,y)\mapsto t=xy$ 描述，$t$ 是底上的局部坐标，$t=0$ 描述节点轨迹。

与复几何的联系相当基本：曲面 $\Sigma$ 的微分同胚群 $\mathrm{Diff}(\Sigma)$（若需要，固定一组指定的标记点）的分类空间，同伦等价于同一曲面上复结构的模轨形。这在双曲情形下精确成立（即除去亏格 $0$ 带一两个标记点、以及亏格 $1$ 不带标记点的情形）。若对节点曲面施加稳定性条件——所有不可约分支都应是双曲的——那么 Lefschetz 纤维化也有万有分类空间，即代数几何中研究得很多的 Deligne--Mumford 模轨形 $\overline{\mathcal{M}}^n_g$。指标 $g$ 与 $n$ 指亏格与标记点个数（在 TQFT 语境中，可把它们看作很小的边界圆），横线表示这些空间是光滑代数曲线模空间 $\mathcal{M}^n_g$ 的紧化。

轨形 $\overline{\mathcal{M}}^n_g$ 有优美的局部结构：其正规交叉边界除子按曲线获得节点的方式标记——拓扑上想象把曲面上一个简单闭圈掐成一个自交点。各层的组合结构编码了一个精细的代数结构（「循环 operad」，见 [15]），它与 CohFT 的结构密切相关。但在最简单的定义中，一个 CohFT 由一族类
\[
H^*\Bigl(\overline{\mathcal{M}}^{p+q}_g;\operatorname{Hom}\bigl(Z(S^1)^{\otimes p},Z(S^1)^{\otimes q}\bigr)\Bigr)
\]
中的元素确定。这些类满足分解条件，条件描述了它们到 $\overline{\mathcal{M}}^n_g$ 边界层的限制，并在此语境中表达了 §3.1.1 中给出的观念所启发的条件：曲面的复合导致线性映射的复合。

\begin{remark}
$B\mathrm{Diff}(\Sigma)$ 与拓扑型为 $\Sigma$ 的光滑黎曼面模轨形之间的同伦等价并不神秘。我们可以把 $B\mathrm{Diff}$ 实现为曲面上度量的模轨形模去微分同胚。但复结构等同于模去共形伸缩的度量，而伸缩空间可缩。所以度量曲面与共形曲面的模空间等价。真正潜藏于背景中的较难结果是：在双曲情形，$\mathrm{Diff}$ 的各分支可缩，从而 $B\mathrm{Diff}\sim B\pi_0\mathrm{Diff}$。
\end{remark}

上同调场论这一概念，源于想要编码 Kähler（或辛）流形 $X$ 的 Gromov--Witten 不变量的结构——这些不变量计数满足给定关联条件的全纯曲线。忽略模空间 $\overline{\mathcal{M}}$ 上的高阶类所得到的Frobenius 代数，就是著名的量子上同调，即 $H^*(X)$ 上杯积的一个（交换）形变。Kontsevich--Manin 公理包含一些 Gromov--Witten 情形特有的约束，最重要的是关于分次的约束。理解 GW 不变量完整结构的一般策略撞上了一个致命障碍：空间 $\overline{\mathcal{M}}^n_g$ 的上同调至今未知。于是理解 CohFT 的结构类似于研究未知环上的模！我们将在 §3.3–§3.5 讨论的扩张 TQFT，其关键动机之一正是 Kontsevich 的方案：从极小数据出发代数地分类 TQFT，再由之重构 Gromov--Witten 不变量。这一方案仍在发展中，才刚刚开始结果。

该理论的一大成功涉及亏格零的部分——它在几何上足够有趣，因为它在代数簇中计数（有理曲线的树）。注意，带标记点的光滑亏格零曲线空间足够简单（$\mathbb{P}^1$ 中互异点的构形），但它的 Deligne--Mumford 紧化 $\overline{\mathcal{M}}^n_0$ 远非平凡，其边界除子以有趣的模式相交。这个空间的上同调被 Keel [17] 完全理解：他给出了一个显式表示，即由（边界除子的庞加莱对偶）生成的自由 $\Z$-代数，模去线性与二次关系。（这些关系容易通过研究遗忘映射 $\overline{\mathcal{M}}^n_0\to\overline{\mathcal{M}}^4_0\cong\mathbb{P}^1$ 的纤维得到；例如，位于该映射不相交纤维中的除子，其交积为零。）

一个更专门但同样重要的成功遵循了 Givental 的想法，他研究了量子上同调（形变后的 $H^*(X)$）成为半单Frobenius 代数的流形的 Gromov--Witten 理论。$\mathbb{P}^n$ 就是这种情形：通常的上同调环变成 $\C[\omega]/\omega^{n+1}=q$；所有环面流形一般也如此。（理论中的参数 $q$ 用来按次数分离全纯曲线的计数。）基于充分的实验证据，Givental 猜想：若量子上同调Frobenius 代数半单，则所有 GW 不变量——回忆它们计数各种亏格与次数的全纯曲线——由该Frobenius 代数本身（连同分次信息）唯一确定。换言之，把一个「朴素的」半单 TQFT（如 §3.1.4 中那样）扩张为完整的 CohFT 的方式唯一，且这个扩张有显式（递归）公式。该猜想在 [29] 中被证实；但若去掉半单性假设，这个结果似乎没有希望推广，而该假设非常强。（「大多数」靶流形不满足它；只有面对大量有理曲线时，我们才能指望把 $H^*(X)$ 上幂零的杯积形变到足够半单。）Kontsevich 研究扩张 TQFT 的动机之一，正是为一般簇的 Gromov--Witten 理论找到类似的结构性结果。
